\scanpage{17}
% archive foliation: 16
\begin{center}- 94 -\end{center}\hfill n\textsuperscript{o} 268

construction, et c'est donc un homomorphisme (Chap.I, \S 1). On notera que si $A$ est complète, $B$ l'est aussi.

De plus, le diviseur $\overline{D(u)}$ étant rationnel sur $k(f(u))$, il existe un diviseur positif $E$ sur $B\times V$ tel que $\overline{D(u)}=E(f(u))$ (par F-VII$_6$ th.12). Les conditions (i) et (ii) sont alors évidemment satisfaites La conditions (iii) est précisément ce que nous avons démontré dans le lemme et la condition (iv) est celle que nous venons d'obtenir.

Le corollaire 1 - est alors facile à démontrer. Je dis que le noyau de $f$ est égal au noyau de l'homomorphisme
\[
w\longrightarrow\operatorname{Cl}[D(w+v)-D(v).]
\]
En effet, si $D(a+v)-D(v)\sim0$ avec $a\in A$, et $v$ générique sur $k(a)$, alors $f(a+v)=f(v)$, et par conséquent $f(a)=0$. Réciproquement, si $f(a)=0$, alors pour $v$ générique sur $k(a)$, $f(a+v)=f(v)$ et $E(f(a+v))=E(f(v))$, autrement dit $\overline{D(a+v)}=\overline{D(v)}$. Cela donne bien $D(a+v)-D(v)\sim0$, et termine notre discussion du quotient dans le théorème du carré.

4 - \uline{Applications du théorème du carré aux variétés abéliennes.}

Dans ce numéro, nous allons employer le fait que la variété de groupe est complète pour obtenir des résultats plus précis sur l'application $\varphi_X:a\longrightarrow\operatorname{Cl}(X_a-X)$ définie à la fin du numéro précédent.

Nous dirons qu'un diviseur $X$ sur une variété abélienne $A$ est \uline{non dégénéré} si le noyau de $\varphi_X$ est fini. Nous démontrerons au cours du chapitre les résultats suivants : Il existe des diviseurs non dégénérés, et pour de tels diviseurs, l'application $\varphi_X$ est surjective. Il existe un entier $d$ dépendant de $A$ seulement, tel que pour n'importe quel diviseur $Y$ on ait $\varphi_Y=0$ si et seulement si $dY$ est algébriquement équivalent à $0$. Pour définir ensuite la
